\chapter{快速上手}

\section{OpenClaw 简介}

\begin{quote}
\textit{"去壳！去壳！"} — 大概是一只太空龙虾说的
\end{quote}

OpenClaw 是一个 AI 智能体 Gateway 网关，支持 WhatsApp、Telegram、Discord、iMessage 等多种聊天平台。
通过单个 Gateway 网关进程将聊天应用连接到 AI 编程智能体，支持本地或远程部署。

\subsection{工作原理}

\begin{center}
\begin{tikzpicture}[
    node distance=0.8cm and 1.5cm,
    every node/.style={font=\small},
    box/.style={draw=blue!50, fill=blue!5, rounded corners=4pt, minimum width=2.8cm, minimum height=0.7cm, align=center},
    gw/.style={draw=orange!60, fill=orange!10, rounded corners=6pt, minimum width=3.5cm, minimum height=1cm, align=center, font=\bfseries},
    arr/.style={-{Stealth[length=2.5mm]}, thick, draw=gray!60}
]
    % Left: Chat apps
    \node[box] (wa) {WhatsApp};
    \node[box, below=of wa] (tg) {Telegram};
    \node[box, below=of tg] (dc) {Discord};
    \node[box, below=of dc] (im) {iMessage};

    % Center: Gateway
    \node[gw, right=2cm of $(tg)!0.5!(dc)$] (gw) {Gateway\\网关};

    % Right: Outputs
    \node[box, right=2cm of gw, yshift=0.8cm] (agent) {AI Agent};
    \node[box, right=2cm of gw, yshift=-0.8cm] (web) {Web UI};

    % Arrows
    \foreach \src in {wa, tg, dc, im} {
        \draw[arr] (\src) -- (gw);
    }
    \draw[arr] (gw) -- (agent);
    \draw[arr] (gw) -- (web);
\end{tikzpicture}
\end{center}

Gateway 网关是会话、路由和渠道连接的\textbf{唯一事实来源}。

\section{三步安装}

\subsection{第一步：安装 OpenClaw}

\begin{lstlisting}[language=bash]
npm install -g openclaw@latest
\end{lstlisting}

\subsection{第二步：运行新手引导}

\begin{lstlisting}[language=bash]
openclaw onboard --install-daemon
\end{lstlisting}

新手引导向导会自动配置：
\begin{itemize}
  \item 模型/认证（推荐 OAuth 或 API 密钥）
  \item 渠道登录（WhatsApp、Telegram 等）
  \item 守护进程（systemd/launchd 开机自启）
\end{itemize}

\subsection{第三步：启动 Gateway}

\begin{lstlisting}[language=bash]
openclaw channels login
openclaw gateway --port 18789
\end{lstlisting}

打开浏览器访问 \texttt{http://127.0.0.1:18789/} 即可开始聊天。

\section{最快体验：Control UI}

无需任何渠道配置，直接用浏览器聊天：

\begin{lstlisting}[language=bash]
openclaw dashboard
\end{lstlisting}

在浏览器中打开 \texttt{http://127.0.0.1:18789/}，即可与智能体对话。

\section{认证配置}

\subsection{Anthropic API 密钥（推荐）}

\begin{lstlisting}[language=bash]
export ANTHROPIC_API_KEY="..."
openclaw models status
\end{lstlisting}

或将密钥写入守护进程配置：

\begin{lstlisting}[language=bash]
cat >> ~/.openclaw/.env <<'EOF'
ANTHROPIC_API_KEY=...
EOF
\end{lstlisting}

\subsection{验证安装}

\begin{lstlisting}[language=bash]
openclaw models status
openclaw doctor
\end{lstlisting}

\section{核心概念速览}

OpenClaw 由以下核心组件构成，后续章节将逐一详解：

\begin{center}
\begin{tabular}{lll}
\toprule
\textbf{概念} & \textbf{说明} & \textbf{配置路径} \\
\midrule
Gateway & AI 网关核心，WS/HTTP 服务器 & \texttt{gateway.*} \\
Agent & AI 助手实例，绑定模型/技能 & \texttt{agents.list[]} \\
Channel & 聊天平台适配器 & \texttt{channels.<name>.*} \\
Binding & 用户/群组到 Agent 的路由 & \texttt{agents.list[].bindings[]} \\
Session & 一次会话，含上下文历史 & \texttt{session.*} \\
Plugin & 扩展模块，注册工具/渠道/Hook & \texttt{plugins.entries.*} \\
Skill & Agent 技能（工具+系统提示） & \texttt{skills.entries.*} \\
Memory & 长期记忆，向量+全文搜索 & \texttt{memorySearch.*} \\
Sandbox & Docker 沙箱隔离执行 & \texttt{agents.defaults.sandbox.*} \\
\bottomrule
\end{tabular}
\end{center}

\section{配置文件}

配置文件位于 \texttt{\textasciitilde{}/.openclaw/openclaw.json}（JSON5 格式）。

\begin{lstlisting}[language=json]
{
  channels: {
    whatsapp: {
      allowFrom: ["+15555550123"],
      groups: { "*": { requireMention: true } }
    }
  },
  messages: {
    groupChat: { mentionPatterns: ["@openclaw"] }
  }
}
\end{lstlisting}

如果你\textbf{不做任何修改}，OpenClaw 将使用内置的 Pi 二进制文件以 RPC 模式运行，并按发送者创建独立会话。
